\section{Эволюция архитектурного подхода}

На раннем этапе исследования естественным вариантом было обучение системы напрямую предсказывать текстовую команду (Steps.txt) по двум состояниям модели – до и после элементарной операции. Однако эксперименты показали, что тип операции определяется достаточно надёжно, а точные численные параметры – значительно хуже. Это связано с чувствительностью параметров к масштабу модели, системе координат и дискретизации поверхности.

В качестве промежуточного решения была предложена модульная архитектура, где после извлечения признаков из двух STL выполнялась классификация типа операции, а затем специализированный блок регрессии предсказывал параметры. Однако точные числа по-прежнему трудно восстановить из дискретизированной геометрии.

Следующий шаг заключался в переосмыслении роли нейросети: вместо предсказания команд, она должна отвечать за распознавание геометрических паттернов и локализацию следов операций. Параметрическая реконструкция переносится на сторону CAD-ядра, где доступны точные геометрические вычисления. Такой подход, получивший название \textit{гибридная сегментационно-классификационная схема}, позволяет:
\begin{itemize}
    \item избежать прямой регрессии численных параметров;
    \item получить интерпретируемый результат в виде «раскраски» модели;
    \item итеративно применять и проверять гипотезы на стороне CAD;
    \item использовать специализированные сети для разных операций.
\end{itemize}

В таблице~\ref{tab:approach-comparison} приведено сравнение рассмотренных подходов.

\begin{table}[h]
\centering
\caption{Сравнение подходов к реализации}
\label{tab:approach-comparison}
\begin{tabular}{|p{4cm}|p{4cm}|p{4cm}|p{3.5cm}|}
\hline
\textbf{Подход} & \textbf{Преимущества} & \textbf{Ограничения} & \textbf{Итоговая оценка} \\
\hline
Прямое предсказание Steps.txt & Явный формат выхода; логичная интеграция через API & Высокая чувствительность к численным параметрам, системе координат и накоплению ошибки & Хорошая исследовательская гипотеза, но недостаточно устойчива \\
\hline
Модульная классификация + регрессия & Разделение типа операции и параметров; расширяемость по операциям & Параметры остаются труднорегрессируемыми из дискретизированной геометрии & Ценный промежуточный шаг \\
\hline
Сегментация полной геометрии & Соответствует природе 3D-данных; даёт интерпретируемый результат; позволяет перенести точные вычисления в CAD & Требует автоматической разметки и балансировки датасета по операциям & Наиболее перспективный подход на текущем этапе \\
\hline
Специализированные сети по операциям & Упрощение задачи; независимое дообучение; повышение управляемости качества & Нужна оркестрация результатов и отдельные выборки под каждый класс & Рекомендуемое направление развития \\
\hline
\end{tabular}
\end{table}

